\documentclass[conference]{IEEEtran}
\IEEEoverridecommandlockouts
% The preceding line is only needed to identify funding in the first footnote. If that is unneeded, please comment it out.
\usepackage{cite}
\usepackage{hyperref}
\usepackage{amsmath,amssymb,amsfonts}
\usepackage{algorithmic}
\usepackage{graphicx}
\usepackage{textcomp}
\usepackage{xcolor}
\def\BibTeX{{\rm B\kern-.05em{\sc i\kern-.025em b}\kern-.08em
    T\kern-.1667em\lower.7ex\hbox{E}\kern-.125emX}}
\begin{document}

\title{Blockchain solutions for incentivization of market-linked upskilling in labor markets}

\author{\IEEEauthorblockN{1\textsuperscript{st} Thomas Scholtz}
	\IEEEauthorblockA{\textit{Department of Computer Science} \\
		\textit{University of Stellenbosch}\\
		Stellenbosch, South Africa \\
		21681147@sun.ac.za}
	\and
	\IEEEauthorblockN{2\textsuperscript{nd} Brink Van Der Merwe}
	\IEEEauthorblockA{\textit{Department of Computer Science} \\
		\textit{University of Stellenbosch}\\
		Stellenbosch, South Africa \\
		abvdm@sun.ac.za}

}

\maketitle

\begin{abstract}
	Blockchain technology has proved to be useful for basic applications such as peer-to-peer electronic cash systems i.e., Bitcoin~\cite{nakamoto2008bitcoin}.  Since this milestone, the technology has evolved to support advanced applications which require more elaborate feature sets. Ethereum~\cite{buterin2014ethereum} is an example of one such form of blockchain technology that allows public deployment of capture-resistant, permissionless contracts, resulting in a platform on which interested parties may attempt to design systems that better address systemic social challenges.

	This paper proposes blockchain solutions for incentivization of market-linked upskilling in South African labor markets with intent to address persistent unemployment in Stellenbosch, South Africa - a problem that is hypothesized to be due in part to the cold-start problem of underprivileged individuals in a dual-speed economy.

	The incentive scheme comprises three blockchain solutions: incentive tokens disbursed to learners on completion of course milestones, verified credentials assigned to learners for completion of courses, and HyperCerts awarded to educators for the effective upskilling of learners via courses.
	The solution(s) is investigated in distinct parts and as a whole by: architecting the components of the solution; simulating the emergent scheme as a sum of the solutions; and surveying relevant stakeholders on the perceived efficacy of such a solution with respect to resolving local unemployment in the town of Stellenbosch.

\end{abstract}

\begin{IEEEkeywords}
	blockchain, tokenomics, HyperCert, Ethereum, incentivisation, unemployment
\end{IEEEkeywords}

\section{Introduction}
Blockchain technology has proved valuable for applications ranging from peer-to-peer electronic cash systems such as Bitcoin~\cite{nakamoto2008bitcoin} to general-purpose computing platforms such as Ethereum~\cite{buterin2014ethereum}. Ethereum offers a Turing-complete scripting environment that enables developers to encode arbitrary rules for ownership, transaction formats, and state transitions. This potential has given rise to a large and growing class of decentralized applications (DApps), with recent surveys cataloguing thousands of deployed DApps across domains including finance, gaming, governance, and identity~\cite{zheng2023dapp}. The emergence of a composable platform with fundamental properties of decentralization, persistence, and auditability, poses a strong case for the design of systems that address pressing systemic social challenges.

Realising ambitious application-layer systems on Ethereum has historically been constrained by the base layer's limited throughput and high transaction costs. A substantial body of recent work addresses these constraints through Layer 2 (L2) scaling architectures including state channels, sidechains, and, most prominently, optimistic and zero-knowledge rollups~\cite{sguanci2021layer2,thibault2022rollups}. Rollup-based L2s now offer orders-of-magnitude higher throughput while inheriting the security guarantees of the underlying L1, rendering them sufficiently mature to reliably host production-scale social and economic applications. It is recognized that are there still pressing technical challenges to be overcome by the Ethereum network (cite); however, the existing utility offered by the network is sufficient for the purposes of this research. The maturity of the network is confirmed by an accelerated interest in blockchain for social good. Sansone et al ~\cite{sansone2023blockchain} identify reliability, transparency, decentralization, and accessibility as the properties through which blockchain can support social innovation, while earlier work documents proof-of-concept deployments spanning humanitarian aid, supply-chain transparency, digital identity, and philanthropy~\cite{marino2018blockchain}. Much of this work nevertheless remains exploratory, with limited engagement with specific, deeply rooted socio-economic problems in a particular context.

South Africa presents one such context. Youth unemployment has persisted above forty-five percent for more than a decade and, in Stellenbosch in particular, is driven in significant part by a dual-speed economy in which underprivileged entrants face a structural cold-start problem: they cannot acquire market-relevant experience without initial access to the labour market, and cannot access the labour market without prior experience or at least highly relevant skills (mostly gated by experience in labour markets) ~\cite{bhorat2016youth,graham2015youth,francis2019inequality}.
Supply-side interventions targeting general employment skills alone have been shown to be insufficient; they must be coupled with demand-side incentives that connect learners, educators, and employers ~\cite{mckenzie2017almp,bhorat2016youth}. It is reasonable to suggest that sophisticated incentives for employers to inform the upskilling of the unemployed may create that appropriate dynamic that is needed for impactful unemployment intervention. There are three benefits in such a dynamic that are immediately apparent: the unemployed learn skills that are highly relevant for employment; the employers are able to 'pre-train' interested individuals in the unemployed population; the employers are rewarded with financial incentive. However, the cold-start problem where unemployed individuals lack the means to provide for themselves in the short-term still remains. The proposed solution to this is the disbursement of incentive tokens on completion of course milestones, which can be used to purchase basic goods at a network of pre-approved vendors, who exchange the tokens for ZAR stablecoin, and then ZAR. The ZAR stablecoins are either funded by government municipalities, or by Corporate Social Responsibility programs of the employer-educators.

While the merit of this hypothesis is subject to further social studies contextualised to Stellenbosch, the focus of this paper is on the technical feasibility of such a system, making the large assumption that the hypothesis is true in the context of this research.

This paper recognises the support for designing application-layer systems on Ethereum and the sufficient need for such systems in the South African labour market context and proposes a blockchain-based incentivization scheme for market-linked upskilling in the South African labour market, composed of three interrelated solutions: fungible incentive tokens disbursed to learners on the completion of course milestones; non-transferable, or soulbound, credential tokens assigned to learners for course completion, in accordance to the decentralized society proposal of Weyl et al.~\cite{weyl2022decentralized}; and HyperCerts~\cite{dalrymple2022hypercerts} awarded to educators for the effective upskilling of learners. The combined ecosystem is designed using frameworks that treat incentive design, governance, and tokenomics as interdependent pillars of a token economy~\cite{oliveira2018token,kivilo2026designing,sadykhov2023detect}. The remainder of the paper architects each component, simulates the resulting system to study its tokenomics and tune parameters, and surveys stakeholders on the perceived efficacy of such a scheme in addressing structural unemployment.

\section{Ease of Use}

\subsection{Maintaining the Integrity of the Specifications}

The IEEEtran class file is used to format your paper and style the text. All margins,
column widths, line spaces, and text fonts are prescribed; please do not
alter them. You may note peculiarities. For example, the head margin
measures proportionately more than is customary. This measurement
and others are deliberate, using specifications that anticipate your paper
as one part of the entire proceedings, and not as an independent document.
Please do not revise any of the current designations.

\section{Prepare Your Paper Before Styling}
Before you begin to format your paper, first write and save the content as a
separate text file. Complete all content and organizational editing before
formatting. Please note sections \ref{AA}--\ref{SCM} below for more information on
proofreading, spelling and grammar.

Keep your text and graphic files separate until after the text has been
formatted and styled. Do not number text heads---{\LaTeX} will do that
for you.

\subsection{Abbreviations and Acronyms}\label{AA}
Define abbreviations and acronyms the first time they are used in the text,
even after they have been defined in the abstract. Abbreviations such as
IEEE, SI, MKS, CGS, ac, dc, and rms do not have to be defined. Do not use
abbreviations in the title or heads unless they are unavoidable.

\subsection{Units}
\begin{itemize}
	\item Use either SI (MKS) or CGS as primary units. (SI units are encouraged.) English units may be used as secondary units (in parentheses). An exception would be the use of English units as identifiers in trade, such as ``3.5-inch disk drive''.
	\item Avoid combining SI and CGS units, such as current in amperes and magnetic field in oersteds. This often leads to confusion because equations do not balance dimensionally. If you must use mixed units, clearly state the units for each quantity that you use in an equation.
	\item Do not mix complete spellings and abbreviations of units: ``Wb/m\textsuperscript{2}'' or ``webers per square meter'', not ``webers/m\textsuperscript{2}''. Spell out units when they appear in text: ``. . . a few henries'', not ``. . . a few H''.
	\item Use a zero before decimal points: ``0.25'', not ``.25''. Use ``cm\textsuperscript{3}'', not ``cc''.)
\end{itemize}

\subsection{Equations}
Number equations consecutively. To make your
equations more compact, you may use the solidus (~/~), the exp function, or
appropriate exponents. Italicize Roman symbols for quantities and variables,
but not Greek symbols. Use a long dash rather than a hyphen for a minus
sign. Punctuate equations with commas or periods when they are part of a
sentence, as in:
\begin{equation}
	a+b=\gamma\label{eq}
\end{equation}

Be sure that the
symbols in your equation have been defined before or immediately following
the equation. Use ``\eqref{eq}'', not ``Eq.~\eqref{eq}'' or ``equation \eqref{eq}'', except at
the beginning of a sentence: ``Equation \eqref{eq} is . . .''

\subsection{\LaTeX-Specific Advice}

Please use ``soft'' (e.g., \verb|\eqref{Eq}|) cross references instead
of ``hard'' references (e.g., \verb|(1)|). That will make it possible
to combine sections, add equations, or change the order of figures or
citations without having to go through the file line by line.

Please don't use the \verb|{eqnarray}| equation environment. Use
\verb|{align}| or \verb|{IEEEeqnarray}| instead. The \verb|{eqnarray}|
environment leaves unsightly spaces around relation symbols.

Please note that the \verb|{subequations}| environment in {\LaTeX}
will increment the main equation counter even when there are no
equation numbers displayed. If you forget that, you might write an
article in which the equation numbers skip from (17) to (20), causing
the copy editors to wonder if you've discovered a new method of
counting.

	{\BibTeX} does not work by magic. It doesn't get the bibliographic
data from thin air but from .bib files. If you use {\BibTeX} to produce a
bibliography you must send the .bib files.

	{\LaTeX} can't read your mind. If you assign the same label to a
subsubsection and a table, you might find that Table I has been cross
referenced as Table IV-B3.

{\LaTeX} does not have precognitive abilities. If you put a
\verb|\label| command before the command that updates the counter it's
supposed to be using, the label will pick up the last counter to be
cross referenced instead. In particular, a \verb|\label| command
should not go before the caption of a figure or a table.

Do not use \verb|\nonumber| inside the \verb|{array}| environment. It
will not stop equation numbers inside \verb|{array}| (there won't be
any anyway) and it might stop a wanted equation number in the
surrounding equation.

\subsection{Some Common Mistakes}\label{SCM}
\begin{itemize}
	\item The word ``data'' is plural, not singular.
	\item The subscript for the permeability of vacuum $\mu_{0}$, and other common scientific constants, is zero with subscript formatting, not a lowercase letter ``o''.
	\item In American English, commas, semicolons, periods, question and exclamation marks are located within quotation marks only when a complete thought or name is cited, such as a title or full quotation. When quotation marks are used, instead of a bold or italic typeface, to highlight a word or phrase, punctuation should appear outside of the quotation marks. A parenthetical phrase or statement at the end of a sentence is punctuated outside of the closing parenthesis (like this). (A parenthetical sentence is punctuated within the parentheses.)
	\item A graph within a graph is an ``inset'', not an ``insert''. The word alternatively is preferred to the word ``alternately'' (unless you really mean something that alternates).
	\item Do not use the word ``essentially'' to mean ``approximately'' or ``effectively''.
	\item In your paper title, if the words ``that uses'' can accurately replace the word ``using'', capitalize the ``u''; if not, keep using lower-cased.
	\item Be aware of the different meanings of the homophones ``affect'' and ``effect'', ``complement'' and ``compliment'', ``discreet'' and ``discrete'', ``principal'' and ``principle''.
	\item Do not confuse ``imply'' and ``infer''.
	\item The prefix ``non'' is not a word; it should be joined to the word it modifies, usually without a hyphen.
	\item There is no period after the ``et'' in the Latin abbreviation ``et al.''.
	\item The abbreviation ``i.e.'' means ``that is'', and the abbreviation ``e.g.'' means ``for example''.
\end{itemize}
An excellent style manual for science writers is \cite{b7}.

\subsection{Authors and Affiliations}
\textbf{The class file is designed for, but not limited to, six authors.} A
minimum of one author is required for all conference articles. Author names
should be listed starting from left to right and then moving down to the
next line. This is the author sequence that will be used in future citations
and by indexing services. Names should not be listed in columns nor group by
affiliation. Please keep your affiliations as succinct as possible (for
example, do not differentiate among departments of the same organization).

\subsection{Identify the Headings}
Headings, or heads, are organizational devices that guide the reader through
your paper. There are two types: component heads and text heads.

Component heads identify the different components of your paper and are not
topically subordinate to each other. Examples include Acknowledgments and
References and, for these, the correct style to use is ``Heading 5''. Use
``figure caption'' for your Figure captions, and ``table head'' for your
table title. Run-in heads, such as ``Abstract'', will require you to apply a
style (in this case, italic) in addition to the style provided by the drop
down menu to differentiate the head from the text.

Text heads organize the topics on a relational, hierarchical basis. For
example, the paper title is the primary text head because all subsequent
material relates and elaborates on this one topic. If there are two or more
sub-topics, the next level head (uppercase Roman numerals) should be used
and, conversely, if there are not at least two sub-topics, then no subheads
should be introduced.

\subsection{Figures and Tables}
\paragraph{Positioning Figures and Tables} Place figures and tables at the top and
bottom of columns. Avoid placing them in the middle of columns. Large
figures and tables may span across both columns. Figure captions should be
below the figures; table heads should appear above the tables. Insert
figures and tables after they are cited in the text. Use the abbreviation
``Fig.~\ref{fig}'', even at the beginning of a sentence.

\begin{table}[htbp]
	\caption{Table Type Styles}
	\begin{center}
		\begin{tabular}{|c|c|c|c|}
			\hline
			\textbf{Table} & \multicolumn{3}{|c|}{\textbf{Table Column Head}}                                                         \\
			\cline{2-4}
			\textbf{Head}  & \textbf{\textit{Table column subhead}}           & \textbf{\textit{Subhead}} & \textbf{\textit{Subhead}} \\
			\hline
			copy           & More table copy$^{\mathrm{a}}$                   &                           &                           \\
			\hline
			\multicolumn{4}{l}{$^{\mathrm{a}}$Sample of a Table footnote.}
		\end{tabular}
		\label{tab1}
	\end{center}
\end{table}

\begin{figure}[htbp]
	\centerline{\includegraphics[width=\columnwidth]{fig1.png}}
	\caption{Example of a figure caption.}
	\label{fig}
\end{figure}

Figure Labels: Use 8 point Times New Roman for Figure labels. Use words
rather than symbols or abbreviations when writing Figure axis labels to
avoid confusing the reader. As an example, write the quantity
``Magnetization'', or ``Magnetization, M'', not just ``M''. If including
units in the label, present them within parentheses. Do not label axes only
with units. In the example, write ``Magnetization (A/m)'' or ``Magnetization
\{A[m(1)]\}'', not just ``A/m''. Do not label axes with a ratio of
quantities and units. For example, write ``Temperature (K)'', not
``Temperature/K''.

\section*{Acknowledgment}

The preferred spelling of the word ``acknowledgment'' in America is without
an ``e'' after the ``g''. Avoid the stilted expression ``one of us (R. B.
G.) thanks $\ldots$''. Instead, try ``R. B. G. thanks$\ldots$''. Put sponsor
acknowledgments in the unnumbered footnote on the first page.
\begin{thebibliography}{00}
	\bibitem{nakamoto2008bitcoin} S. Nakamoto, ``Bitcoin: A peer-to-peer electronic cash system,'' 2008. [Online]. Available: https://bitcoin.org/bitcoin.pdf

	\bibitem{buterin2014ethereum} V. Buterin, ``A next-generation smart contract and decentralized application platform,'' Ethereum White Paper, 2014. [Online]. Available: https://ethereum.org/en/whitepaper/

	\bibitem{zheng2023dapp} P. Zheng, Z. Jiang, J. Wu, and Z. Zheng, ``Blockchain-based decentralized application: A survey,'' \emph{IEEE Open Journal of the Computer Society}, vol.~4, pp.~121--133, 2023, doi: 10.1109/OJCS.2023.3251854.

	\bibitem{sguanci2021layer2} C. Sguanci, R. Spatafora, and A. M. Vergani, ``Layer 2 blockchain scaling: A survey,'' 2021, arXiv:2107.10881. [Online]. Available: https://arxiv.org/abs/2107.10881

	\bibitem{thibault2022rollups} L. T. Thibault, T. Sarry, and A. S. Hafid, ``Blockchain scaling using rollups: A comprehensive survey,'' \emph{IEEE Access}, vol.~10, pp.~93039--93054, 2022, doi: 10.1109/ACCESS.2022.3200051.

	\bibitem{sansone2023blockchain} G. Sansone, D. Leone, and F. Schiavone, ``Blockchain for social good and stakeholder engagement: Evidence from a case study,'' \emph{Corporate Social Responsibility and Environmental Management}, vol.~30, no.~5, pp.~2182--2192, 2023, doi: 10.1002/csr.2477.

	\bibitem{marino2018blockchain} F. Marino and A. Juels, ``Blockchain for social good,'' in \emph{Proc. 4th EAI Int. Conf. Smart Objects and Technologies for Social Good (GoodTechs)}, Bologna, Italy, 2018, pp.~166--168, doi: 10.1145/3284869.3284881.

	\bibitem{bhorat2016youth} H. Bhorat, A. Cassim, and D. Tseng, ``The state of youth unemployment in South Africa,'' Brookings Institution, Washington, DC, Policy Brief, 2016.

	\bibitem{graham2015youth} L. Graham and C. Mlatsheni, ``Youth unemployment in South Africa: Understanding the challenge and working on solutions,'' in \emph{South African Child Gauge 2015}, A. De Lannoy, S. Swartz, L. Lake, and C. Smith, Eds. Cape Town, South Africa: Children's Institute, University of Cape Town, 2015, pp.~51--59.

	\bibitem{francis2019inequality} D. Francis and E. Webster, ``Poverty and inequality in South Africa: Critical reflections,'' \emph{Development Southern Africa}, vol.~36, no.~6, pp.~788--802, 2019, doi: 10.1080/0376835X.2019.1666703.

	\bibitem{weyl2022decentralized} E. G. Weyl, P. Ohlhaver, and V. Buterin, ``Decentralized society: Finding Web3's soul,'' SSRN, Working Paper 4105763, May 2022. [Online]. Available: https://papers.ssrn.com/sol3/papers.cfm?abstract\_id=4105763

	\bibitem{dalrymple2022hypercerts} D. A. Dalrymple, H. Holtebrinck, T. Kuhn, J. Miller, R. L. Ocampo, and P. Schmidt, ``Hypercerts: A new primitive for public goods funding,'' Protocol Labs, Whitepaper, 2022. [Online]. Available: https://www.protocol.ai/blog/hypercert-new-primitive/

	\bibitem{oliveira2018token} L. Oliveira, L. Zavolokina, I. Bauer, and G. Schwabe, ``To token or not to token: Tools for understanding blockchain tokens,'' in \emph{Proc. 39th Int. Conf. Information Systems (ICIS)}, San Francisco, CA, USA, Dec. 2018.

	\bibitem{kivilo2026designing} S. Kivilo, A. Norta, M. Hattingh, S. Avanzo, and L. Pennella, ``Designing a token economy: Incentives, governance, and tokenomics,'' 2026, arXiv:2602.09608. [Online]. Available: https://arxiv.org/abs/2602.09608

	\bibitem{sadykhov2023detect} A. Sadykhov, G. Goodell, D. Dhinakaran, P. Horshalom, and P. Treleaven, ``Decentralized token economy theory (DeTEcT): Token pricing, stability and governance for token economies,'' \emph{Frontiers in Blockchain}, vol.~6, art.~1298330, 2023, doi: 10.3389/fbloc.2023.1298330.

	\bibitem{mckenzie2017almp} D. McKenzie, ``How effective are active labor market policies in developing countries? A critical review of recent evidence,'' \emph{The World Bank Research Observer}, vol.~32, no.~2, pp.~127--154, 2017, doi: 10.1093/wbro/lkx001.

\end{thebibliography}
\vspace{12pt}

\end{document}

